%% kapa: exam=1 round=34 year=2023 period=1 form=A subject=civil_law subject_ko=민법 seq=21 number=21
%% source: Q-Net 감정평가사 기출문제 (cst00309) artlSeq=5211789 file=raw/1008/5211789/2023년도 제34회 감정평가사 1차시험 1교시.pdf
%% conversion: kordoc 4.24.51 -> tools/gen1.py (auto); review=proofread-sonnet; flags=none
\begin{question}{21}
물권의 객체에 관한 설명으로 옳은 것은? (다툼이 있으면 판례에 따름)
\begin{choices}
  \item 지상권은 물건이 아니므로 저당권의 객체가 될 수 없다.
  \item 법률상 공시방법이 인정되지 않는 유동집합물이라도 특정성이 있으면 이를 양도담보의 목적으로 할 수 있다.
  \item 저당권과 질권은 서로 다른 물권이므로 하나의 물건에 관하여 동시에 성립할 수 있다.
  \item 토지소유권은 토지의 상하에 미치므로 지상공간의 일부만을 대상으로 하는 구분지상권은 원칙적으로 허용되지 않는다.
  \item 기술적인 착오 없이 작성된 지적도에서의 경계가 현실의 경계와 다르다면, 토지소유권의 범위는 원칙적으로 현실의 경계를 기준으로 확정하여야 한다.
\end{choices}
\end{question}
\answer{2}
